% ============================================================
%  صفحه‌ی عنوان و معرفی
% ============================================================
\begin{titlepage}
\thispagestyle{empty}

% ─── نوار بالایی ──────────────────────────────────────────────
\begin{tikzpicture}[remember picture, overlay]
  % پس‌زمینه‌ی بالا
  \fill[PrimaryDeep] (current page.north west) rectangle
       ([yshift=-4.5cm]current page.north east);
  % نوار طلایی
  \fill[AccentGold] ([yshift=-4.5cm]current page.north west) rectangle
       ([yshift=-4.9cm]current page.north east);
  % نوار پایین
  \fill[PrimaryLight] ([yshift=-4.9cm]current page.north west) rectangle
       ([yshift=-5.3cm]current page.north east);

  % متن بالا (راست‌چین فارسی)
  \node[anchor=east, text=white, font=\Large\bfseries]
    at ([xshift=-2cm, yshift=-2cm]current page.north east)
    {\rl{مقاله‌ی تحلیلی–تطبیقی در فلسفه‌ی سیاسی}}; % FIXED: \rl

  \node[anchor=east, text=NeutralLight, font=\normalsize]
    at ([xshift=-2cm, yshift=-3.5cm]current page.north east)
    {\rl{بررسی تاریخی، مکتبی و تطبیقی مفهوم آزادی در اندیشه‌ی سیاسی}}; % FIXED: \rl

  % نوار پایین صفحه
  \fill[PrimaryDeep] (current page.south west) rectangle
       ([yshift=2cm]current page.south east);
  \fill[AccentGold] ([yshift=2cm]current page.south west) rectangle
       ([yshift=2.3cm]current page.south east);

  \node[anchor=east, text=white, font=\small]
    at ([xshift=-2cm, yshift=1cm]current page.south east)
    {{\bfseries \rl{مهدی سالم}} \quad\textbar\quad ۳ اسفند ۱۴۰۴ \quad\textbar\quad ۲۲ فوریه ۲۰۲۶}; % FIXED: \rl
\end{tikzpicture}

\vspace{5cm}

% ─── عنوان اصلی ───────────────────────────────────────────────
\begin{center}
{\fontsize{42}{50}\selectfont\bfseries\color{PrimaryDeep}آزادی}

\vspace{0.3cm}
{\fontsize{28}{36}\selectfont\color{NeutralDark} به مثابه\ldots}

\vspace{0.8cm}
\begin{tikzpicture}
  \draw[line width=2pt, AccentGold] (0,0) -- (12,0);
  \fill[AccentGold] (6,0) circle (4pt);
\end{tikzpicture}

\vspace{0.8cm}
{\Large\color{TextSecond}
  جغرافیای مفهومی، امواج تاریخی و تحلیل تطبیقی\\[0.3cm]
  گفتمان‌های رقیب در فلسفه‌ی سیاسی
}

\vspace{1.5cm}

% ─── اطلاعات نویسنده ──────────────────────────────────────────
\begin{tcolorbox}[
  enhanced, width=10cm,
  colback=BgCool, colframe=PrimaryMid,
  rounded corners, drop shadow,
]
\centering
{\large\bfseries\color{PrimaryDeep} مهدی سالم}\\[0.3cm]
{\normalsize\color{TextSecond}
  پژوهشگر و نویسنده در حوزه‌ی\\
  گذار دموکراتیک و فلسفه‌ی سیاسی
}\\[0.4cm]
\small\color{NeutralDark}
\lr{\url{MahdiSalem.com}}\\[0.2cm] % FIXED: \lr{\url{...}}
{\color{AccentGold}\rule{6cm}{0.5pt}}\\[0.2cm]
{\small ۳ اسفند ۱۴۰۴ \textbar{} ۲۲ فوریه ۲۰۲۶}
\end{tcolorbox}

\end{center}

\vspace{1cm}

% ─── خلاصه‌ی اجرایی ───────────────────────────────────────────
\begin{tcolorbox}[
  enhanced, breakable,
  colback=BgWarm, colframe=AccentGold,
  title={\bfseries چکیده},
  fonttitle=\large\bfseries,
  attach boxed title to top right={yshift=-2mm},
  boxed title style={colback=AccentGold, colframe=AccentAmber},
]
این مقاله‌ی تحلیلی–تطبیقی، مفهوم آزادی را از سه منظر بررسی می‌کند: نخست، \textbf{جغرافیای مفهومی} آن را در قالب نقشه‌ای از تعبیرسازی‌های تاریخی ترسیم می‌کند؛ دوم، \textbf{امواج تحول} این مفهوم را از دوران باستان تا دوران دیجیتال دنبال می‌کند؛ و سوم، \textbf{تحلیل تطبیقی} مکاتب رقیب—از لیبرالیسم تا مارکسیسم، از جمهوری‌خواهی تا فمینیسم، از سنت تحلیلی تا سنت قاره‌ای—را با تمرکز بر دشمنان مفهومی، انسان‌شناسی‌های زیربنایی، و تنش‌های ساختاری ارائه می‌دهد. در پیوست‌ها، سنت اسلامی–ایرانی با عمق بیشتری تحلیل می‌شود.

\vspace{0.3cm}
\small\color{NeutralDark}
\textbf{کلیدواژه‌ها:} آزادی منفی، آزادی مثبت، عدم سلطه، قابلیت، خودفرمانی، رهایی، سنت جمهوری‌خواه، لیبرالیسم، فلسفه‌ی سیاسی تطبیقی
\end{tcolorbox}

\end{titlepage}

% ─── فهرست مطالب ──────────────────────────────────────────────
\newpage
\thispagestyle{empty}
\begin{center}
{\Large\bfseries\color{PrimaryDeep} فهرست مطالب}
\vspace{0.5cm}
\end{center}
\tableofcontents
\newpage
